\chapter{Convenciones y Patrones de Desarrollo}

\section{Patrones de C\'odigo}

\subsection{Server Actions}

Todas las server actions siguen estas convenciones:

\begin{lstlisting}[style=typescript, caption=Patron de Server Action]
'use server';                          // 1. Directiva

import { db, schema } from '@vetrinaria/db';
import { eq } from 'drizzle-orm';
import { revalidatePath } from 'next/cache';
import { getCurrentClinicId } from '@/lib/clinic';

export async function getItems() {     // 2. Prefijo get para lecturas
  const clinicId = await getCurrentClinicId();  // 3. Contexto clinica
  return db                             // 4. Query Drizzle
    .select()
    .from(schema.items)
    .where(eq(schema.items.clinicId, clinicId))
    .orderBy(schema.items.name);
}

export async function createItem(data) { // 5. Prefijo create/update/delete
  const clinicId = await getCurrentClinicId();
  await db.insert(schema.items).values({ ... });
  revalidatePath('/items');             // 6. Revalidacion de cache
}
\end{lstlisting}

\subsection{Tipos de Datos}

\begin{itemize}[leftmargin=*]
    \item \textbf{Con Zod:} Cuando hay validaci\'on de entrada significativa
          (clientes, cl\'{\i}nicas). Usar \texttt{z.infer} para tipos.
    \item \textbf{Sin Zod:} Para casos simples donde el tipado de TypeScript
          es suficiente (pacientes, citas, facturas).
    \item \textbf{Enums:} Definidos como \texttt{const as const} arrays y
          usados como \texttt{text()} con restricci\'on de enum en Drizzle.
\end{itemize}

\subsection{Manejo de Enums en Drizzle}

Los enums de PostgreSQL se definen como arrays de TypeScript y se referencian
en el esquema:

\begin{lstlisting}[style=typescript, caption=Patron de enums en Drizzle]
export const appointmentStatusEnum = [
  'scheduled', 'confirmed', 'checked_in',
  'in_exam', 'checked_out', 'cancelled', 'no_show',
] as const;

status: text('status', { enum: appointmentStatusEnum })
  .notNull().default('scheduled')
\end{lstlisting}

Los cast con \texttt{as any} son necesarios al insertar strings en campos
con restricci\'on de enum.

\subsection{JSONB Tipado}

Los campos JSONB usan \texttt{\$type<T>()} para tipado en TypeScript:

\begin{lstlisting}[style=typescript, caption=Patron JSONB tipado]
weight: jsonb('weight')
  .$type<{ value: number; date: string }[]>()
  .default([]),
\end{lstlisting}

\subsection{\'Indices}

Cada tabla tiene \'{\i}ndices para \texttt{clinic\_id} y las columnas
m\'as consultadas. Las tablas con consultas compuestas tienen \'{\i}ndices
compuestos (ej: \texttt{appointments\_date\_status\_idx}).

\section{Convenciones de Nomenclatura}

\subsection{Archivos y Directorios}

\begin{longtable}{|l|l|}
\hline
\textbf{Elemento} & \textbf{Convenci\'on} \\
\hline
Directorios de p\'aginas & kebab-case (\texttt{medical-records}) \\
Archivos de esquemas DB & snake\_case (\texttt{medical\_records.ts}) \\
Archivos de componentes & kebab-case (\texttt{appointment-form.tsx}) \\
Archivos de server actions & kebab-case (\texttt{medical-records.ts}) \\
Directorios de API routes & kebab-case con [params] \\
\hline
\end{longtable}

\subsection{C\'odigo}

\begin{longtable}{|l|l|}
\hline
\textbf{Elemento} & \textbf{Convenci\'on} \\
\hline
Tablas Drizzle & snake\_case en DB, camelCase en TypeScript \\
Columnas Drizzle & snake\_case en DB, camelCase en TypeScript \\
Funciones & camelCase \\
Tipos/interfaces & PascalCase \\
Constantes compartidas & UPPER\_CASE \\
Props de componentes & PascalCase \\
\hline
\end{longtable}

\section{Internacionalizaci\'on}

\begin{itemize}[leftmargin=*]
    \item UI en espa\~nol (etiquetas, mensajes, documentaci\'on).
    \item Enums en ingl\'es en base de datos, con traducciones en
          \texttt{@vetrinaria/shared/constants}.
    \item Datos semilla con nombres y direcciones en espa\~nol.
\end{itemize}

\section{Estructura de P\'aginas}

Cada m\'odulo de negocio sigue esta estructura:

\begin{lstlisting}[style=bash, caption=Estructura tipica de modulo de negocio]
app/(dashboard)/
+-- modulo/
|   +-- page.tsx          # Listado
|   +-- new/
|   |   +-- page.tsx      # Creacion
|   +-- [id]/
|       +-- page.tsx      # Detalle
|       +-- edit/
|           +-- page.tsx  # Edicion
+-- lib/actions/
    +-- modulo.ts          # Server actions
+-- components/
    +-- modulo-form.tsx   # Formulario reutilizable
\end{lstlisting}

\section{Seguridad}

\begin{itemize}[leftmargin=*]
    \item \textbf{Server-only:} M\'odulos sensibles (backup, clinic context)
          marcados con \texttt{'server-only'}.
    \item \textbf{Contrase\~nas:} Hasheadas con bcryptjs (12 rondas).
    \item \textbf{Autenticaci\'on:} NextAuth con middleware que protege
          todas las rutas del dashboard.
    \item \textbf{Multi-tenant:} Filtrado por \texttt{clinic\_id} en todas
          las consultas.
    \item \textbf{Credenciales de prueba:} Incluidas en la p\'agina de login
          para facilitar desarrollo.
\end{itemize}

\section{Rendimiento}

\begin{itemize}[leftmargin=*]
    \item \textbf{\'Indices:} 38 \'{\i}ndices B-tree en PostgreSQL.
    \item \textbf{Consultas paralelas:} \texttt{Promise.all} para consultas
          independientes (dashboard, m\'ultiples fetches).
    \item \textbf{Caching:} \texttt{cache()} de React para deduplicaci\'on
          de datos del servidor.
    \item \textbf{Server Components:} M\'aximo renderizado en servidor,
          m\'{\i}nimo JavaScript de cliente.
    \item \textbf{PNPM:} Instalaci\'on eficiente de dependencias con
          content-addressable store.
    \item \textbf{Turborepo:} Cache de build para evitar recomputaciones.
\end{itemize}
